\documentclass[10pt,twoside]{article}
\usepackage[T1]{fontenc}
\usepackage[utf8]{inputenc}
\usepackage{amsmath,amssymb,amsthm,amsfonts,mathrsfs,mathtools,bm}
\usepackage{graphicx,xcolor,tikz,booktabs,multirow,tabularx,array,float}
\usepackage{caption,subcaption,enumitem,setspace,verbatim}
\usepackage{natbib,hyperref,fancyhdr,geometry,xspace,pdflscape}
\usepackage{algorithm,algpseudocode}
\geometry{papersize={16.5cm,23.5cm},top=2cm,bottom=1.5cm,left=1.5cm,right=1.5cm}
\hypersetup{colorlinks=true,linkcolor=blue,citecolor=magenta,urlcolor=cyan}
\pagestyle{fancy}
\fancyfoot{}
\renewcommand{\headrulewidth}{.01cm}
\fancyhead[RE,LO]{\thepage}
\newtheorem{dfn}{Definition}[section]
\newtheorem{thm}[dfn]{Theorem}
\newtheorem{pro}[dfn]{Proposition}
\newtheorem{rem}[dfn]{Remark}
\newtheorem{lem}[dfn]{Lemma}
\newtheorem{cor}[dfn]{Corollary}
\newtheorem{eg}[dfn]{Example}
\newtheorem{alg}{Algorithm}
\let\oldtitle\title
\renewcommand{\bibname}{References}
\numberwithin{equation}{section}
\graphicspath{{./}{main/p.1096/}{main/p.1096/Images/}{Images/}}
\makeatletter
\def\@maketitle{%
  \begin{center}
  \vspace*{1cm}
  {\Large\bfseries \@title\par}
  \vspace{0.8em}
  {\bfseries \@author\par}
  \vspace{0.5em}
  \end{center}
  \rule{\textwidth}{0.2mm}\\
}
\makeatother
\begin{document}
\thispagestyle{empty}

\begin{quote}


\includegraphics[width=5.98453in,height=0.99667in]{main//p.1096//vertopal_02dcbf921eaa45acad2f5a464a87959f/media/image1.jpeg}

\end{quote}



\section{Application of Best Linear Unbiased Prediction (BLUP) in

Stability Analysis and Genotype × Environment Interaction

Studies}



\textbf{Naser Sabaghnia}\footnote{\begin{quote}

  \textbf{Corresponding Author:} sabaghnia@maragheh.ac.ir

  \end{quote}}



\begin{quote}

\textbf{University of Maragheh, Iran}

\end{quote}



\textbf{Abstract:}



Genotype × environment interaction (GEI) is a major challenge in plant

breeding, as it complicates the identification of genotypes that are

both high-performing and stable across diverse environments.

Multi-environment trials (METs) generate complex datasets, often

unbalanced and incomplete, where traditional ANOVA or regression-based

stability metrics are limited. Best Linear Unbiased Prediction (BLUP),

within a mixed model framework, provides a robust approach for

estimating genotypic performance while accounting for random effects,

environmental variation, and heterogeneous residuals. BLUP enables

shrinkage-based predictions, improving accuracy over raw means or BLUE,

particularly in unbalanced datasets. Furthermore, incorporation of

pedigree or marker-derived relationship matrices allows exploitation of

genetic correlations, enhancing prediction of untested genotypes and

facilitating genomic selection. BLUP-based methods can also quantify

stability through variance of predicted genotypic values across

environments, factor-analytic decomposition of G×E, and

reliability-adjusted indices, providing actionable insights for

selection. This manuscript reviews theoretical foundations, practical

implementation, and computational approaches for BLUP in METs, including

hierarchical and factorial designs. It also discusses integration with

factor-analytic models for complex GEI, the use of combined pedigree and

marker information, and strategies for handling missing data. BLUP-based

stability analysis represents a flexible and powerful tool for modern

plant breeding, enabling selection of widely adapted and specifically

adapted genotypes while maximizing genetic gain under diverse

environmental conditions.



\textbf{Keywords:} Mixed models, Plant breeding, Multi-environment

trials



\textbf{Mathematics Subject Classification (2010):} 62P10, 62J12, 62H25,

92D10.



\section{Introduction}



Plant breeders aim to select genotypes that combine high yield potential

with stability across diverse environments. Phenotypic performance of

genotypes is determined by their genetic makeup (G), environmental

conditions (E), and the interaction between genotype and environment

(G×E). While classical analysis of variance (ANOVA) partitions total

phenotypic variation into G, E, and GEI components, it does not provide

precise predictions of genotypic performance, especially in unbalanced

trials or when data from multiple environments and years are available

(Smith et al., 2001). Stability analysis is often employed to quantify

the consistency of genotype performance across environments. Classical

methods include the regression-based Finlay-Wilkinson approach, AMMI

(Additive Main Effects and Multiplicative Interaction) models, and

Eberhart-Russell parameters (Lynch \& Walsh, 1998). However, these

approaches have limitations in handling unbalanced data, missing

observations, and complex pedigree structures.



Mixed models and Best Linear Unbiased Prediction (BLUP) have emerged as

powerful tools for addressing these limitations (Piepho \& Möhring,

2006). BLUP provides estimates of genotypic effects by explicitly

modeling random effects and variance-covariance structures,

accommodating unbalanced datasets, repeated measurements, and

heterogeneous variances. Furthermore, BLUP can incorporate pedigree or

marker-based genetic relationships, improving prediction accuracy for

selection in breeding programs. This manuscript reviews the theoretical

foundations of BLUP in the context of GEI and stability analysis,

discusses model formulation for hierarchical and factorial experimental

designs, and presents practical applications and considerations for

plant breeding programs.



\section{Stability Experiments}



Multi-environment trials are used to evaluate the performance of

genotypes across diverse environmental conditions. Typical designs

include randomized complete block designs (RCBD), incomplete block

designs, lattice designs, and factorial experiments. Trials are often

conducted over multiple years and locations to capture environmental

variability. Phenotypic data collected from METs are subject to both

spatial and temporal variation (Fig. 1). The statistical model for

phenotypic data is generally expressed as:



\[y_{ijk} = \mu + g_{i} + e_{j} + (ge)_{ij} + \varepsilon_{ijk}

\]



BLUP is derived from the general linear mixed model (LMM):



\[\mathbf{y} = \mathbf{X}\mathbf{\beta} + \mathbf{Zu} + \mathbf{\varepsilon}\mathbf{

}\]



\[\mathbf{u} \sim N(0,\mathbf{G}),\mathbf{\varepsilon} \sim N(0,\mathbf{R})\]



where \(\mathbf{G}\) represents the variance--covariance structure among

random effects, which can be simple (homogeneous) or complex

(heterogeneous or structured by pedigree or markers). \(\mathbf{R}\)

represents residual covariance, which may also accommodate spatial or

temporal correlation (Piepho \& Möhring, 2006).



The BLUP of \(\mathbf{u}\) is obtained by solving the mixed model

equations:



\[\begin{bmatrix}

\mathbf{X}^{'}\mathbf{R}^{- 1}\mathbf{X} & \mathbf{X}^{'}\mathbf{R}^{- 1}\mathbf{Z} \\

\mathbf{Z}^{'}\mathbf{R}^{- 1}\mathbf{X} & \mathbf{Z}^{'}\mathbf{R}^{- 1}\mathbf{Z} + \mathbf{G}^{- 1} \\

\end{bmatrix}\begin{bmatrix}

\widehat{\mathbf{\beta}} \\

\widehat{\mathbf{u}} \\

\end{bmatrix} = \begin{bmatrix}

\mathbf{X}^{'}\mathbf{R}^{- 1}\mathbf{y} \\

\mathbf{Z}^{'}\mathbf{R}^{- 1}\mathbf{y} \\

\end{bmatrix}

\]



Genetic relationships among genotypes can be included via the numerator

relationship matrix (\(\mathbf{A}\)) or marker-based similarity matrices

(\(\mathbf{S}\)). The variance-covariance structure of genotypic effects

can be expressed as:



\[\mathbf{G} = \sigma_{g}^{2}\mathbf{A}\text{or}\mathbf{G} = \sigma_{g}^{2}\mathbf{S}\mathbf{

}\]



\[\mathbf{G} = r_{A}^{2}\mathbf{A} + r_{S}^{2}\mathbf{S}\mathbf{

}\]



\begin{longtable}[]{@{}

  >{\raggedright\arraybackslash}p{(\columnwidth - 0\tabcolsep) * \real{1.0000}}@{}}

\toprule()

\begin{minipage}[b]{\linewidth}\raggedright

\includegraphics[width=5.32in,height=1.86314in]{main//p.1096//vertopal_02dcbf921eaa45acad2f5a464a87959f/media/image2.jpg}

\end{minipage} \\

\midrule()

\endhead

Figure 1. Decomposition of phenotypic variance \\

\bottomrule()

\end{longtable}



BLUP can be applied to hierarchical (nested) and factorial designs:

Hierarchical (nested), so genotypes nested within families. Model allows

estimation of within- and between-family effects. Factorial (crossed),

like diallel crosses or hybrids from two heterotic pools. Model can

include general combining ability (GCA), specific combining ability

(SCA), and test-cross effects.



For example, a diallel model with GCA and SCA effects is expressed as:



\[y_{ijk} = \mu + g_{i} + g_{j} + s_{ij} + e_{k} + \varepsilon_{ijk}

\]



Stability is assessed by decomposing GEI effects or predicting genotype

performance across environments with low fluctuations (Fig. 2).

BLUP-based approaches included (i) BLUP of genotype-by-environment

effects which predicts genotypic performance in specific environments

(Schabenberger \& Gotway, 2005). (ii) Reliability-based stability

measures so every stability variance can be estimated using predicted

genotypic values. (iii) Factor-analytic models, which approximate G×E

covariance structure using latent factors; BLUPs of factor scores

provide stability estimates.



\begin{longtable}[]{@{}

  >{\raggedright\arraybackslash}p{(\columnwidth - 0\tabcolsep) * \real{1.0000}}@{}}

\toprule()

\begin{minipage}[b]{\linewidth}\raggedright

\includegraphics[width=4.54478in,height=1.99222in]{main//p.1096//vertopal_02dcbf921eaa45acad2f5a464a87959f/media/image3.jpeg}

\end{minipage} \\

\midrule()

\endhead

Figure 2. Effect of shrinkage on raw means \\

\bottomrule()

\end{longtable}



Simulation can illustrate BLUP performance. For example, an experiment

with 50 genotypes tested in 6 environments, with 3 replicates per

environment, can be simulated with heterogeneous genetic variances. BLUP

provides shrinkage-adjusted predictions with higher correlation to true

genotypic values than raw means or BLUEs, particularly in unbalanced

datasets.



\section{SAS codes}



Regarding accept the assumptions of Model 1, the effect of the general

adaptability ability (GAA) of genotype i is equal to u+v\emph{i}. The

deviation from the sum u+v\emph{i} represents the specific adaptability

ability effect (SAA\emph{\textsubscript{ik}}) of genotype i in

environment k. This effect consists of a linear component (the effect of

environment dk) and a nonlinear component (the interaction effect

(vd)\textsubscript{ik}). To determine the effects of

GAA\textsubscript{i} and SAA\textsubscript{ik}, the mean phenotypic

values across replications are calculated. In this way, the results of

testing n genotypes across m environments can be arranged in an n×m

table, where: x\textsubscript{i} is the sum of all phenotypic values for

genotype i, x\textsubscript{k} is the total phenotypic sum in

environment k, X.. is the grand total of all phenotypic values,

x\textsubscript{ik} is the phenotypic value of genotype i in environment

k. The GAA and SAA effects are calculated using the following formulas

(Fig. 1):



\[x_{i.} = \sum_{k}^{}x_{ik} = x_{i1} + x_{i2} + \ldots + x_{im}\]



\[x_{nk} = \sum_{k}^{}x_{ik} = x_{1k} + x_{2k} + \ldots + x_{nk}\]



\[x_{ik} = u + v_{i} + d_{k} + {(vd)}_{ik}\]



\[x_{ik} = u + {GAA}_{i} + {SAA}_{ik}\]



\begin{longtable}[]{@{}

  >{\raggedright\arraybackslash}p{(\columnwidth - 0\tabcolsep) * \real{1.0000}}@{}}

\toprule()

\begin{minipage}[b]{\linewidth}\raggedright

\includegraphics[width=4.23313in,height=1.83843in]{main//p.1096//vertopal_02dcbf921eaa45acad2f5a464a87959f/media/image4.jpeg}

\end{minipage} \\

\midrule()

\endhead

Figure 1. Schematic illustration of stability evaluation \\

\bottomrule()

\end{longtable}



The variance of random deviations for the mean phenotypic value in a

specific environment x\textsubscript{ik} is computed as:



\[{var}_{xik} = \frac{\sigma_{e}^{2}}{c} = \sigma^{2}\]



The variance of the difference between the means of two genotypes is

determined as:



\[{var}_{(x_{ik}/x_{nm})} = 2\sigma^{2}\]



To obtain the standard error of the differences between effects, the

square root of the corresponding variance must be taken. The least

significant difference (LSD) is then obtained by multiplying the

standard error by a constant magnitude. Comparison of genotypes in terms

of general adaptability can be performed by comparing

GAA\textsubscript{i} values. To estimate the interaction strength of

genotype i with environments, the genotype × environment interaction

variance can be used. To measure the stability of genotype i, it is

recommended to use the variance of SAA\textsubscript{i}. Therefore,

stability can be defined as the capacity of a genotype to maintain its

phenotypic expression under various environmental conditions through

organized mechanisms.



\section{Stability analysis}



Practical code examples for fitting BLUP in multi-environment trials

(METs), including stability analysis using factor-analytic models are:



proc mixed data=MET\_data;



class Genotype Environment Replicate;



model Yield = Environment / solution;



random Genotype Genotype*Environment;



/* Genotype and GxE effects as random */



lsmeans Genotype / pdiff cl;



run;



Genotype and Genotype*Environment are treated as random effects to

obtain BLUPs.



Environment is fixed to capture environmental means.



lsmeans Genotype provides predicted genotypic values across

environments.



Practical code examples for fitting BLUP in multi-environment trials

(METs), including



Factor-analytic model for G×E are:



proc mixed data=MET\_data method=REML;



class Genotype Environment;



model Yield = Environment;



random Genotype / type=fa(1) g;



repeated Environment / subject=Genotype type=ar(1);



ods output SolutionR=BLUPs;



run;



type=fa(1) specifies a factor-analytic covariance structure for G×E.



repeated allows modeling residual correlation across environments.



BLUPs output contains predicted genotypic values adjusted for G×E.



\begin{longtable}[]{@{}

  >{\raggedright\arraybackslash}p{(\columnwidth - 0\tabcolsep) * \real{1.0000}}@{}}

\toprule()

\begin{minipage}[b]{\linewidth}\raggedright

\includegraphics[width=3.74772in,height=2.97015in]{main//p.1096//vertopal_02dcbf921eaa45acad2f5a464a87959f/media/image5.jpeg}

\end{minipage} \\

\midrule()

\endhead

Figure 2. Flowchart of breeding and selection strategy \\

\bottomrule()

\end{longtable}



\section{Breeding for GAA while considering

stability}



The BLUP accounts for missing data, unbalanced replication, and

heterogeneous variance, yielding more accurate estimates of genotypic

performance across environments compared to classical ANOVA or

regression-based stability analysis. In factorial and diallel crosses,

BLUP allows estimation of GCA and SCA effects even for untested hybrids

(Bernardo, 1994). BLUP predictions enable the derivation of stability

metrics across environments, so variance of predicted genotype

performance across environments reflects stability. Also,

factor-analytic models separate main effects from interaction patterns,

highlighting genotypes with stable performance (low sensitivity to

environmental variation) or specific adaptation (high performance in

certain environments).



BLUP can model complex covariance structures among environments, such as

autoregressive or spatially correlated residuals (Edwards \& Jannink,

2006). Genotypic effects can be predicted for untested environments,

enabling selection for wide or targeted adaptation. Marker information

can enhance BLUP predictions by improving estimation of genetic

relationships, especially for genotypes with incomplete pedigree data.

Combined models of A and S matrices allow optimal utilization of both

pedigree and marker data.



BLUP provides a flexible, robust framework for analyzing MET data and

assessing GEI. Compared with classical stability metrics, BLUP improves

prediction of genotypic performance under unbalanced and incomplete data

(Fig. 3). Also, it accounts for complex variance-covariance structures,

including spatial, temporal, and pedigree effects. BLUP Supports

hierarchical and factorial designs, enabling dissection of GCA, SCA, and

within-family variation. Finally, it Facilitates integration of pedigree

and marker-based genetic relationships for genomic prediction.



The main assumptions of BLUP are normally distributed random effects and

correctly specified variance-covariance structures. In practice, REML

estimation is robust to moderate deviations from normality, although

outliers may affect inference (Piepho \& Möhring, 2006). Accurate

variance component estimation is crucial, particularly in plant

breeding, where the number of genotypes is often limited. Bayesian or

shrinkage approaches can help address heterogeneity among crosses or

environments (Edwards \& Jannink, 2006). BLUP-based stability analysis

also allows flexible interpretation of GEI (Schabenberger \& Gotway,

2005). Factor-analytic models and prediction-based stability metrics

provide information on general adaptability and specific adaptation,

guiding breeders in selecting genotypes for wide or targeted

environmental adaptation.



\begin{longtable}[]{@{}

  >{\raggedright\arraybackslash}p{(\columnwidth - 0\tabcolsep) * \real{1.0000}}@{}}

\toprule()

\begin{minipage}[b]{\linewidth}\raggedright

\includegraphics[width=5.4403in,height=1.52166in]{main//p.1096//vertopal_02dcbf921eaa45acad2f5a464a87959f/media/image6.jpg}

\end{minipage} \\

\midrule()

\endhead

Figure 3. Workflow for BLUP-based stability analysis \\

\bottomrule()

\end{longtable}



\section{Conclusion}



BLUP is a powerful tool for stability analysis and GEI studies in plant

breeding. By incorporating random effects, variance--covariance

structures, and pedigree or marker information, BLUP enables precise

prediction of genotypic performance across environments. Its ability to

handle unbalanced and incomplete data, coupled with flexible model

specification, makes it superior to classical methods in

multi-environment trials. The integration of BLUP with genomic

information and hierarchical experimental designs further enhances

selection efficiency. Future research should focus on robust variance

estimation in small datasets, robustness to non-normality, and practical

applications in modern breeding programs.



\section{References}



Bernardo, R. (1994). Predicting hybrid performance from parental lines:

Implications for the design of breeding programs. \emph{Theoretical and

Applied Genetics,} 89, 917--924.



Piepho, H.P., Möhring, J. (2006). Computing heritability and selection

response from unbalanced plant breeding trials. \emph{Genetics,} 173,

1885--1894.



Smith, A.B., Cullis, B.R., Thompson, R. (2001). Analyzing variety by

environment data using multiplicative mixed models and adjustments for

spatial field trend. \emph{Biometrics,} 57, 1138--1147.



Lynch, M., Walsh, B. (1998). Genetics and Analysis of Quantitative

Traits. Sinauer Associates. MA, USA.



Schabenberger, O., Gotway, C.A. (2005). \emph{Statistical Methods for

Spatial Data Analysis.} CRC Press. FL, USA



Edwards, J., Jannink, J.L. (2006). Bayesian modeling of

genotype-by-environment interactions for assessing crop stability.

\emph{Theoretical and Applied Genetics,} 112, 1247--1255.

\end{document}
